%DS10_Trang
%Đề cũ số 10. Đề thường: chỉ đáp án; câu khó có hướng dẫn ngắn.

\section{ĐỀ SỐ \thesection }
\setcounter{bt}{0}

%% Dạng 1. Thực hiện phép tính — 4 ý: a cơ bản, b tính hợp lý, c và d nhiều ngoặc và lũy thừa.
\begin{bt} Thực hiện phép tính:
	\begin{listEX}
		\item $46+(-28)-13$;
		\item $(-15)\cdot 24+(-15)\cdot 16$;
		\item $96:\left[4^{2}-2^{3}\right]+(-1)^{2025}$;
		\item $(-3)^{2}\cdot 4-\left\{50-\left[2^{5}-7\cdot(-2)^{0}\right]\right\}$.
	\end{listEX}
	\loigiai{a) $5$ \hfill b) $-600$ \hfill c) $11$ \hfill d) $11$.}
\end{bt}

%% Dạng 2. Tìm x — 4 ý: a cơ bản, b trị tuyệt đối, c ngoặc lồng, d có x ở cả hai vế.
\begin{bt} Tìm số nguyên $x$, biết:
	\begin{listEX}
		\item $7x=-63$;
		\item $|4x-2|=18$;
		\item $3\left\{2\left[x-(-2)\right]-6\right\}=18$;
		\item $8-3(x-2)=2(x+1)-x$.
	\end{listEX}
	\loigiai{a) $x=-9$ \hfill b) $x=5$ hoặc $x=-4$ \hfill c) $x=4$ \hfill d) $x=3$.}
\end{bt}

%% Dạng 3. Tập hợp — 2 ý, số nguyên và số tự nhiên.
\begin{bt} Biết số $0$ là số tự nhiên. Gọi $A$ là tập hợp các số nguyên $x$ thỏa mãn $-9\le x\le 9$, $x$ chia hết cho $3$ và $x$ không chia hết cho $6$. Gọi $B$ là tập hợp các số tự nhiên nhỏ hơn $20$, chia hết cho $3$ và không chia hết cho $2$.
	\begin{enumerate}
		\item Liệt kê các phần tử của $A$ và của $B$.
		\item Liệt kê các số vừa thuộc $A$ vừa thuộc $B$. Liệt kê các số thuộc $A$ nhưng không thuộc $B$.
	\end{enumerate}
	\loigiai{a) $A=\{-9;\,-3;\,3;\,9\}$, $B=\{3;\,9;\,15\}$ \hfill b) Phần chung: $\{3;\,9\}$; thuộc $A$ nhưng không thuộc $B$: $\{-9;\,-3\}$.}
\end{bt}

%% Dạng 4. Ước và bội — thực tế cắt thanh gỗ, học sinh tự nhận ra ƯCLN.
\begin{bt} Người ta có những thanh gỗ dài $48$ cm và những thanh gỗ dài $36$ cm. Người ta cắt các thanh này thành những đoạn bằng nhau, dài nhất có thể, và không dư gỗ.
	\begin{enumerate}
		\item Mỗi đoạn dài bao nhiêu xăng-ti-mét?
		\item Mỗi thanh $48$ cm cắt được bao nhiêu đoạn? Mỗi thanh $36$ cm cắt được bao nhiêu đoạn?
	\end{enumerate}
	\loigiai{a) $12$ cm \hfill b) $4$ đoạn và $3$ đoạn.}
\end{bt}

%% Dạng 5. Toán thực tế — vận tốc, quãng đường, thời gian.
\begin{bt} Quãng đường từ nhà bạn Lan đến thư viện dài $2\,400$ m. Vận tốc đi bộ của bạn Lan là $80$ m/phút. Vận tốc đi xe đạp gấp $3$ lần vận tốc đi bộ.
	\begin{enumerate}
		\item Bạn Lan đi bộ hết bao nhiêu phút?
		\item Bạn Lan đi xe đạp hết bao nhiêu phút?
		\item Đi xe đạp ít hơn đi bộ bao nhiêu phút?
	\end{enumerate}
	\loigiai{a) $30$ phút \hfill b) $10$ phút \hfill c) $20$ phút.}
\end{bt}

%% Dạng 6. Hình học thực tế — lối đi bao quanh sân.
\begin{bt}
	\immini[thm]{
		Một mảnh vườn hình chữ nhật dài $25$ m, rộng $16$ m. Xung quanh vườn, sát bốn cạnh, có một lối đi rộng $2$ m. Phần bên trong lối đi dùng để trồng hoa.
		\begin{enumerate}
			\item Tính diện tích phần trồng hoa.
			\item Tính diện tích lối đi.
			\item Biết giá lát lối đi là $200\,000$ đồng/m$^{2}$. Tính số tiền lát lối đi.
		\end{enumerate}
	}{
		\begin{tikzpicture}[scale=0.22,line width=0.8pt,line join=round,>=stealth]
			\draw (0,0) rectangle (25,16);
			\draw (2,2) rectangle (23,14);
			\node at (12.5,8) {Trồng hoa};
			\draw[<->] (0,-1.4)--(25,-1.4) node[midway,below]{$25$ m};
			\draw[<->] (26.2,0)--(26.2,16) node[midway,right]{$16$ m};
			\draw[<->] (-1.2,14)--(-1.2,16) node[midway,left]{$2$ m};
		\end{tikzpicture}
	}
	\loigiai{a) $252$ m$^{2}$ \hfill b) $148$ m$^{2}$ \hfill c) $29\,600\,000$ đồng.}
\end{bt}

%% Dạng 7. Nâng cao — rút gọn thương của hai tổng lũy thừa cùng cơ số.
\begin{bt} Tính nhanh:
	\[
	\left(2^{7}+2^{8}+2^{9}\right):\left(2^{5}+2^{6}+2^{7}\right).
	\]
	\loigiai{Hướng dẫn: đặt thừa số chung $2^{7}$ ở tử và $2^{5}$ ở mẫu.\\
		Kết quả bằng $4$.}
\end{bt}
